\documentclass[10pt,a4paper]{article}
\usepackage{fontspec}
\setmainfont{Georgia}
\newfontfamily\georgiahead{Georgia}
\usepackage[spanish,es-tabla]{babel}
\usepackage[left=1.6cm, right=1.6cm, top=1.6cm, bottom=1.6cm, footskip=0.7cm]{geometry}
\usepackage{graphicx}
\usepackage{amsmath,amssymb,amsfonts}
\usepackage{xcolor}
\usepackage{fancyhdr}
\usepackage{booktabs}
\usepackage{tabularx}
\usepackage{array}
\usepackage{float}
\usepackage{multicol}
\usepackage{tcolorbox}
\tcbuselibrary{skins, breakable}
\usepackage{hyperref}
\usepackage{titlesec}
\usepackage{caption}
\usepackage{setspace}
\usepackage{microtype}
\usepackage{tocloft}

% Paleta de colores sobria
\definecolor{gsnavy}{HTML}{0D233A}
\definecolor{gsblue}{HTML}{1E3A8A}
\definecolor{riskred}{HTML}{B91C1C}
\definecolor{aluargreen}{HTML}{00843D}
\definecolor{aluargold}{HTML}{D97706}
\definecolor{gsgray}{HTML}{4A4A4A}
\definecolor{lightgray}{HTML}{F8FAFC}
\definecolor{bordergray}{HTML}{D1D5DB}

\hypersetup{
    colorlinks=true,
    linkcolor=gsnavy,
    citecolor=gsnavy,
    urlcolor=gsblue,
    pdftitle={Informe de Valuación Aluar S.A.I.C. - UNCuyo},
}

% Formato de títulos
\titleformat{\section}{\georgiahead\large\bfseries\color{gsnavy}}{\thesection}{0.8em}{}[\color{gsnavy}\titlerule]
\titleformat{\subsection}{\georgiahead\normalsize\bfseries\color{gsblue}}{\thesubsection}{0.6em}{}
\titleformat{\subsubsection}{\georgiahead\small\bfseries\color{gsgray}}{\thesubsubsection}{0.5em}{}
\titlespacing*{\section}{0pt}{6pt}{2pt}
\titlespacing*{\subsection}{0pt}{4pt}{2pt}
\titlespacing*{\subsubsection}{0pt}{3pt}{1pt}

\newcommand{\coverheading}[1]{{\georgiahead\bfseries\color{gsnavy}\normalsize #1}\par\vspace{1pt}}

% Configuración de Captions
\captionsetup{labelfont={bf,color=gsnavy}, textfont={small,color=gsgray}, skip=3pt}

% Configuración de TOC
\renewcommand{\cftsecleader}{\cftdotfill{\cftdotsep}}
\renewcommand{\cftsecfont}{\scriptsize\bfseries\color{gsnavy}}
\renewcommand{\cftsecpagefont}{\scriptsize\bfseries\color{gsnavy}}
\renewcommand{\cftsubsecfont}{\tiny\color{gsgray}}
\renewcommand{\cftsubsecpagefont}{\tiny\color{gsgray}}
\setlength{\cftbeforesecskip}{1pt}
\setlength{\cftbeforesubsecskip}{0.5pt}

% Encabezados y Pie de Página
\pagestyle{fancy}
\fancyhf{}
\setlength{\headheight}{14pt}
\lhead{\small \textbf{ALUAR S.A.I.C. (ALUA.BA)} $\cdot$ Equity Research}
\rhead{\small FCE $\cdot$ UNCuyo}
\lfoot{\color{gsgray}\scriptsize \textsf{Analista: Federico Agustín Chillón | federico.chillon@fce.uncu.edu.ar}}
\rfoot{\color{gsgray}\scriptsize \textsf{Página \thepage}}
\renewcommand{\headrulewidth}{0.6pt}
\renewcommand{\footrulewidth}{0.4pt}

\setlength{\parindent}{1.2em}
\setlength{\floatsep}{3pt plus 1pt minus 1pt}
\setlength{\textfloatsep}{4pt plus 1pt minus 1pt}
\setlength{\intextsep}{3pt plus 1pt minus 1pt}
\setlength{\parskip}{1.5pt plus 0.5pt minus 0.5pt}
\renewcommand{\arraystretch}{0.92}

\begin{document}
\renewcommand{\topfraction}{0.95}
\renewcommand{\bottomfraction}{0.95}
\renewcommand{\floatpagefraction}{0.85}
\renewcommand{\textfraction}{0.05}

% ═════════════════════════════════════════════════════════════════════════════
% HOJA 1: PORTADA
% ═════════════════════════════════════════════════════

\noindent
\begin{minipage}[b]{0.65\textwidth}
    {\fontsize{22}{26}\selectfont \color{gsnavy} \textbf{Aluar Aluminio Argentino}} \\[0.15cm]
    {\large \color{gsgray} \textsf{Informe de Valuación Fundamental y Análisis de Riesgo}}
\end{minipage}%
\begin{minipage}[b]{0.35\textwidth}
    \raggedleft
    {\small \color{gsgray} \textbf{Análisis Financiero}} \\
    {\small \color{gsgray} FCE $\cdot$ UNCuyo} \\
    {\small \color{gsgray} 5 de agosto de 2026}
\end{minipage}

\vspace{0.15cm}
\noindent\textcolor{gsnavy}{\rule{\textwidth}{1.2pt}}
\vspace{0.08cm}

\noindent
{\scriptsize \color{gsgray} \textbf{Aviso legal y descargo de responsabilidad:} Documento de investigación académica y análisis financiero independiente con fines exclusivamente formativos y de divulgación. No constituye oferta de compra, venta o asesoramiento de inversión en los términos de la Ley N.\textsuperscript{o}~26.831. Las valuaciones y conclusiones representan estimaciones del autor fundamentadas en estados financieros públicos auditados y modelos financieros estocásticos.}

\vspace{0.15cm}

\noindent
\begin{minipage}[t]{0.55\textwidth}
    \coverheading{Resumen Ejecutivo}
    {\small El modelo de valuación aplicado a Aluar Aluminio Argentino S.A.I.C. (ALUA.BA) arroja una recomendación técnica de \textbf{COMPRAR} con un precio objetivo fundamental a 12 meses de \textbf{ARS 1.236,00} por acción (Target Integrado \textbf{ARS 1.255,60} al incorporar la opción real de expansión eólica PEAL V), frente a una cotización de mercado de ARS 982,50 (+25,8\% / +27,8\% de retorno esperado). El precio actual descuenta una tasa terminal implícita excesivamente contractiva ($g^* = 0,69\%$), subestimando la ventaja en costos energéticos y el proceso de desapalancamiento operativo proyectado.}

    \vspace{0.15cm}
    \coverheading{Fundamentos Centrales de la Tesis}
    \begin{enumerate}
        \setlength{\itemsep}{2pt}
        \item \textbf{Estructura de Costos de Primer Cuartil:} Matriz energética 96\% autogenerada (78\% renovable vía Futaleufú y PEAL), con costo en efectivo C1 en USD 1.680/t (percentil 29\% global).
        \item \textbf{Generación Dolarizada con Costos Fijos en Pesos:} 80\% de las ventas se exportan a precios fijados en el LME, mientras que los costos operativos domésticos se reducen en términos reales ante ajustes cambiarios.
        \item \textbf{Compresión de Prima Soberana y Desapalancamiento:} La caída del spread EMBI+ (441 pb) sitúa el WACC en 7,06\% ($\lambda=0,20$), y el cese del ciclo intensivo de inversión eólica reduce la deuda neta a $<1,0\text{x}$ EBITDA hacia FY2028E.
    \end{enumerate}
\end{minipage}%
\hfill
\begin{minipage}[t]{0.42\textwidth}
    \vspace{-0.1cm}
    \begin{tcolorbox}[colback=white, colframe=gsnavy, boxrule=0.8pt, arc=3pt, left=4pt, right=4pt, top=3pt, bottom=3pt]
        \begin{center}
            {\scriptsize \color{gsgray} Dictamen del Modelo Fundamental} \\[0.02cm]
            {\fontsize{15}{18}\selectfont \sffamily \textbf{\color{aluargreen} COMPRAR}} \\[0.05cm]
            {\scriptsize \color{gsgray} Target Base (Dumrauf): \textbf{ARS 1.236,00} (USD 0,78)} \\
            {\scriptsize \color{gsgray} Opción Real PEAL V: \textbf{+ARS 19,60} (USD 0,01)} \\
            {\small \color{gsnavy} \textbf{Target Integrado: ARS 1.255,60} (USD 0,79)} \\
            {\scriptsize \color{gsgray} Cotización Spot: ARS 982,50 $\mid$ \textbf{Retorno: +27,8\%}}
        \end{center}
        \vspace{0.04cm}
        \color{gsnavy}\rule{\textwidth}{0.4pt}
        \vspace{0.04cm}
        \scriptsize
        \begin{tabular}{@{}l@{\extracolsep{\fill}}r@{}}
            \textbf{Costo de Capital (WACC):} & 7,06\% \\
            \textbf{Costo del Equity (Ke, $\lambda = 0,20$):} & 9,30\% \\
            \textbf{Beta Apalancado (Hamada):} & 0,888 \\
            \textbf{Prima de Mercado (ERP Damodaran):} & 4,18\% \\
            \textbf{Prima País Ponderada ($\lambda\times$EMBI+):} & 0,88\% \\
            \textbf{Costo de la Deuda (Kd post-tax):} & 2,47\% \\
            \textbf{Estructura D/E Objetivo:} & 0,486 \\
            \textbf{Crecimiento Terminal ($g$):} & 2,00\% \\
            \textbf{Potencial Base / Integrado:} & +25,8\% / +27,8\% \\
            \textbf{Tipo de Cambio Implícito (CCL):} & 1.584,25 \\
        \end{tabular}
        \vspace{0.04cm}
        \color{bordergray}\rule{\textwidth}{0.4pt}
        \vspace{0.03cm}
        {\tiny \color{gsgray} \textbf{Nota de calibración.} Valuación canónica según Dumrauf (Cap. 14) con $CRP = \lambda \times EMBI+$ y Opción Real por mínimos cuadrados de Longstaff-Schwartz. Datos al 23-jul-2026.\par}
    \end{tcolorbox}
\end{minipage}

\vspace{0.15cm}
\noindent
\begin{tcolorbox}[colback=white, colframe=gsblue!40, boxrule=0.5pt, top=2pt, bottom=2pt, left=4pt, right=4pt, title={\small\sffamily \textbf{Catalizadores y Factores de Sensibilidad (Horizonte 12--18 Meses)}}, coltitle=gsnavy]
    \footnotesize
    \begin{itemize}
        \setlength{\itemsep}{0pt}
        \item[\color{aluargreen}$\uparrow$] \textbf{Certificación ASI (Aluminio Verde):} Acceso a primas comerciales en la Unión Europea y exención de aranceles de carbono en frontera (CBAM).
        \item[\color{aluargreen}$\uparrow$] \textbf{Normalización Cambiaria y Régimen RIGI:} Reducción de costos de importación de bienes de capital y estabilidad fiscal para futuros proyectos renovables.
        \item[\color{riskred}$\downarrow$] \textbf{Apreciación del Tipo de Cambio Real (TCR):} Incremento de costos operativos en dólares (salarios y servicios locales).
        \item[\color{riskred}$\downarrow$] \textbf{Contracción del Ciclo del LME:} Desaceleración industrial internacional que sitúe la cotización por debajo de USD 2.200/t.
    \end{itemize}
\end{tcolorbox}

\vspace{0.08cm}
\noindent{\scriptsize \textbf{Aviso legal:} Documento de estudio financiero. Las estimaciones son teóricas y no constituyen asesoramiento financiero.}

\clearpage

% ═════════════════════════════════════════════════════════════════════════════
% HOJA 2: ÍNDICE GENERAL
% ═════════════════════════════════════════════

\begin{center}
    {\large\georgiahead\bfseries\color{gsnavy} Índice General de Secciones y Contenidos}
\end{center}
\vspace{-0.2cm}
\hrule height 0.8pt
\vspace{0.2cm}

\begin{multicols}{2}
\small
\makeatletter
\@starttoc{toc}
\makeatother
\end{multicols}

\vspace{0.2cm}
\hrule height 0.4pt
\vspace{0.15cm}

\noindent
{\scriptsize \color{gsgray} \textbf{Estructura del Documento:} Comprende 17 Secciones de Desarrollo, Apéndice Contable con Estados Financieros Auditados por PwC (FY2020--FY2025), Proyecciones Explícitas de Flujo de Fondos (FCFF 2026E--2030E), 31 Figuras de Análisis Financiero y Referencias Bibliográficas.}

\clearpage

% ═════════════════════════════════════════════════════════════════════════════
% SECCIÓN 1: DESCRIPCIÓN DE LA COMPAÑÍA
% ═════════════════════════════════════════════
\section{Descripción de la Compañía: Integración Vertical y Estructura Operativa}

\subsection{Perfil Operativo: Capacidad Instalada e Integración Productiva}
Aluar Aluminio Argentino S.A.I.C. es el único productor primario de aluminio integrado en Argentina y uno de los mayores complejos industriales de Sudamérica. Fundada en 1974, la compañía opera su planta de reducción electrolítica en Puerto Madryn (Provincia del Chubut) con una capacidad instalada nominal de \textbf{460.000 toneladas anuales}. Su portafolio abarca aluminio primario líquido y sólido (lingotes estándar, lingotes T, barrotes de extrusión, alambrón y placas de laminación), junto con una división de semielaborados y elaborados de mayor valor agregado (perfiles extruidos y laminados para construcción, transporte y embalaje). 

Alrededor del \textbf{80\% del volumen físico producido se destina a la exportación} (Estados Unidos, Brasil, Europa y Japón), liquidado en dólares estadounidenses sobre cotizaciones del London Metal Exchange (LME) más primas regionales, lo que dota a la firma de un flujo operativo vinculado a moneda dura.

\begin{figure}[H]
\centering
\begin{minipage}[t]{0.485\textwidth}
    \centering
    \includegraphics[width=\linewidth,height=4.6cm,keepaspectratio]{../03_Modelo_y_Codigo/figuras/figura_01.png}
    \caption{\textbf{Evolución Histórica de Aluar (1974--2026)}: Hitos de expansión de capacidad instalada y transición energética.}
    \label{fig:01}
\end{minipage}\hfill
\begin{minipage}[t]{0.485\textwidth}
    \centering
    \includegraphics[width=\linewidth,height=4.6cm,keepaspectratio]{../03_Modelo_y_Codigo/figuras/figura_02.png}
    \caption{\textbf{Estructura Accionaria}: Concentración del 45,98\% en el grupo de control (Familia Madanes) y 54,02\% flotante/ANSES.}
    \label{fig:02}
\end{minipage}
\end{figure}

\subsection{Estructura Accionaria y Control Corporativo}
Según la Memoria Anual y Estados Financieros al 30-jun-2025, el capital social está integrado por \textbf{2.800 millones de acciones ordinarias} escriturales de \$1 valor nominal y un voto por acción, listadas en BYMA bajo la especie \textbf{ALUA.BA}. Las sociedades vinculadas a la familia Madanes (CIPAL S.A., Estancias San Javier de Catamarca S.A. y Angerona Group Trust) concentran en conjunto el \textbf{45,98\%} de los derechos de voto (\textbf{Figura \ref{fig:02}}), conformando el grupo de control frente al 54,02\% distribuido entre la ANSES-FGS (27,47\%) y el capital flotante en el mercado doméstico e internacional.

\subsection{Ventajas Competitivas: Matriz Energética y Posición de Costos C1}
La fundición de aluminio primario es un proceso electrointensivo donde la energía representa entre el 30\% y 40\% del costo total de producción. Tres factores sustentan el foso defensivo de la compañía:
\begin{itemize}
    \item \textbf{Autonomía Energética Propia (96\% del Consumo):} La planta de Puerto Madryn es abastecida en un 54\% por la Central Hidroeléctrica Futaleufú (320 MW asignados contractualmente bajo concesión), 24\% por el Parque Eólico Aluar (PEAL, 200 MW en operación) y 18\% por ciclos combinados térmicos propios a gas natural. Sólo el 4\% restante se toma de la red nacional (SADI/CAMMESA), blindando a la compañía ante fluctuaciones en las tarifas eléctricas spot.
    \item \textbf{Posición en el Primer Cuartil de Costos Globales:} Con un costo en efectivo C1 de \textbf{USD 1.680 por tonelada} (percentil 29\% de la curva global de costos CRU/Wood Mackenzie), Aluar se ubica entre los productores más eficientes del mundo, manteniendo márgenes operativos positivos en fases contractivas del ciclo internacional.
    \item \textbf{Terminal Portuaria Exclusiva:} Muelle propio de aguas profundas en Puerto Madryn para la descarga directa de bauxita y alúmina importada y embarque de exportaciones con economías de escala en fletes.
\end{itemize}

\subsection{Estrategia Corporativa: Certificación ASI y Mezcla de Elaborados}
Los objetivos estratégicos comprenden: (1) alcanzar el 100\% de abastecimiento renovable mediante la Etapa V del Parque Eólico (PEAL V / La Flecha); (2) obtener la certificación ASI (*Aluminium Stewardship Initiative*) para capturar primas de precio por aluminio verde en Europa; y (3) defender la cuota en productos elaborados. La participación de elaborados en el volumen total se contrajo temporalmente debido a la recesión del mercado interno, pasando de 6,0\% en FY2023 a 3,3\% en FY2025:

\begin{table}[H]
\centering
\caption{\textbf{Participación de la División Elaborados en el Volumen Total de Ventas} {\small\color{gsgray} (Memoria y EEFF Aluar FY2023–FY2025)}}
\small
\begin{tabular}{lccc}
\toprule
\textbf{Métrica de Volumen} & \textbf{FY2023} & \textbf{FY2024} & \textbf{FY2025} \\
\midrule
Ventas Elaborados (t) & 23.104 & 17.766 & 13.498 \\
Ventas Totales (t) & 383.834 & 418.014 & 402.991 \\
\textbf{Participación Elaborados} & \textbf{6,0\%} & \textbf{4,3\%} & \textbf{3,3\%} \\
\bottomrule
\end{tabular}
\end{table}

\subsection{Análisis Estratégico FODA}
\begin{itemize}
    \item \textbf{Fortalezas:} Monopolio primario integrado en Argentina; 96\% de energía autogenerada; costos C1 en primer cuartil mundial; ingresos 80\% en moneda dura.
    \item \textbf{Oportunidades:} Primas comerciales por certificación ASI; beneficios impositivos del RIGI para PEAL V; compresión de la prima de riesgo soberano.
    \item \textbf{Debilidades:} Dependencia de alúmina importada; capacidad nominal operando al 100\% (460 kt); apalancamiento transitorio por inversiones eólicas.
    \item \textbf{Amenazas:} Ciclo internacional del LME; apreciación del tipo de cambio real que encarece costos laborales en USD; excedentes de fundición en Asia.
\end{itemize}

\subsection{Desempeño Financiero Reciente (FY2020--FY2025)}
\begin{table}[H]
\centering
\caption{\textbf{Evolución de Indicadores Financieros Auditados FY2020–FY2025 (USD MM)}}
\small
\begin{tabular}{lcccccc}
\toprule
\textbf{Métrica Financiera} & \textbf{FY2020} & \textbf{FY2021} & \textbf{FY2022} & \textbf{FY2023} & \textbf{FY2024} & \textbf{FY2025} \\
\midrule
Ingresos (Ventas Netas) & 823,5 & 513,5 & 654,7 & 647,8 & 915,9 & 1.092,6 \\
EBITDA & 116,5 & 107,4 & 209,7 & 105,7 & 204,7 & 163,3 \\
Margen EBITDA & 14,2\% & 20,9\% & 32,0\% & 16,3\% & 22,4\% & 14,9\% \\
Deuda Neta & 351,3 & 137,4 & 57,3 & 169,9 & 211,2 & 525,8 \\
Deuda Neta / EBITDA & 3,02x & 1,28x & 0,27x & 1,61x & 1,03x & 3,22x \\
\bottomrule
\end{tabular}
\end{table}

\noindent
{\tiny \textbf{Nota.} Cifras en USD convertidas al CCL de cierre de cada ejercicio. Deuda Neta auditada según estados financieros consolidados.\par}

% ═════════════════════════════════════════════════════════════════════════════
% SECCIÓN 2: CONTEXTO MACROECONÓMICO Y MECÁNICA DE PROYECCIÓN
% ═════════════════════════════════════════════
\section{Contexto Macroeconómico y Mecánica Operativa de Proyección \texorpdfstring{($P \times Q$)}{(P x Q)}}

\subsection{Entorno Macroeconómico y Dinámica del Riesgo Soberano}
Las proyecciones del modelo se apoyan en el escenario de estabilización macroeconómica argentino (\textbf{Figura \ref{fig:03}}), con convergencia inflacionaria y reactivación gradual de la actividad industrial. La compresión del spread de riesgo soberano EMBI+ desde máximos de 2.400 pb en 2022 a 441 pb al cierre de julio de 2026 (\textbf{Figura \ref{fig:04}}) reduce de forma directa la tasa de descuento requerida para el capital.

\begin{figure}[H]
\centering
\begin{minipage}[t]{0.485\textwidth}
    \centering
    \includegraphics[width=\linewidth,height=4.6cm,keepaspectratio]{../03_Modelo_y_Codigo/figuras/figura_03.png}
    \caption{\textbf{Contexto Macroeconómico}: PBI real e inflación proyectada en Argentina (2020--2030E).}
    \label{fig:03}
\end{minipage}\hfill
\begin{minipage}[t]{0.485\textwidth}
    \centering
    \includegraphics[width=\linewidth,height=4.6cm,keepaspectratio]{../03_Modelo_y_Codigo/figuras/figura_04.png}
    \caption{\textbf{Compresión del Riesgo País (EMBI+)}: Descenso del spread soberano desde 2.400 pb a 441 pb.}
    \label{fig:04}
\end{minipage}
\end{figure}

\subsection{Mecánica Operativa de Ingresos: Volumen Saturado y Reversión del LME}
La proyección de ingresos descansa en la descomposición operativa $Ingresos = P \times Q$:
\begin{enumerate}
    \item \textbf{Volumen Físico ($Q$):} La planta opera a plena capacidad (460.000 Tn nominales). Con un factor de utilización histórico del 94\%, el volumen despachado se proyecta constante en:
    \begin{equation}
        Q = 460.000\text{ Tn} \times 0,94 = \mathbf{432.400\text{ Tn/año}}
    \end{equation}
    \item \textbf{Precio Realizado ($P$):} En FY2026E los ingresos alcanzan USD 1.591,4 MM impulsados por el precio spot del LME en el 1T26 ($USD\text{ }3.173,2/\text{t}$). Para el período 2027E--2030E, el precio realizado converge linealmente a su nivel de equilibrio de largo plazo:
    \begin{equation}
        P_{2030} = LME_{\text{rev}} + Premium_{\text{elab}} = 2.587,1 + 645,7 = \mathbf{USD\text{ }3.232,8/\text{t}}
    \end{equation}
    \item \textbf{Calibración del Margen EBITDA ($k$):} Se calibra el parámetro de eficiencia operativa $k$ sobre los datos reales del 1T26 ($Ventas = USD\text{ }416\text{ MM}$, $EBITDA = USD\text{ }106\text{ MM}$, $C1 = USD\text{ }1.680/\text{t}$):
    \begin{equation}
        k = \frac{Margen_{EBITDA, Q1}}{\left(\frac{P_{Q1} - C1}{P_{Q1}}\right)} = \frac{106 / 416}{\left(\frac{3.848,3 - 1.680}{3.848,3}\right)} = \mathbf{0,4521}
    \end{equation}
    El margen EBITDA de cada período se deduce como $Margen_t = k \times \left(\frac{P_t - 1.680}{P_t}\right)$.
    \item \textbf{Variación de Capital de Trabajo ($\Delta$NWC) en FY2026E:} Se proyecta según la mediana histórica del capital de trabajo operativo sobre ventas ($NWC/\text{Ventas} = 49,42\%$):
    \begin{equation}
        \Delta NWC_{2026E} = (Ventas_{2026E} - Ventas_{2025}) \times 49,42\% = (1.591,4 - 1.092,6) \times 0,4942 = \mathbf{USD\text{ }246,5\text{ MM}}
    \end{equation}
    Esta absorción inicial refleja la recomposición de stocks y créditos operativos tras el salto en ventas, dando paso a una liberación recurrente de $+USD\text{ }23,9\text{ MM/año}$ en el resto del horizonte explícito.
\end{enumerate}

\subsection{Curva de Tasas y Convergencia del Riesgo País (Modelo AR(1))}
Para modelar la trayectoria esperada del EMBI+ se estima un proceso autorregresivo $\text{AR}(1)$ mensual:
\begin{equation}
    EMBI_t = \mu (1 - \phi) + \phi EMBI_{t-1} + \epsilon_t, \quad \hat{\phi} = 0,7582, \; \hat{\mu} = 1.614 \text{ pb}
\end{equation}

\begin{figure}[H]
\centering
\begin{minipage}[t]{0.485\textwidth}
    \centering
    \includegraphics[width=\linewidth,height=4.6cm,keepaspectratio]{../03_Modelo_y_Codigo/figuras/figura_05.png}
    \caption{\textbf{Trayectoria del Riesgo País}: Proyección estocástica AR(1) frente al supuesto corporativo base (441 pb).}
    \label{fig:05}
\end{minipage}\hfill
\begin{minipage}[t]{0.485\textwidth}
    \centering
    \includegraphics[width=\linewidth,height=4.6cm,keepaspectratio]{../03_Modelo_y_Codigo/figuras/figura_06.png}
    \caption{\textbf{Aluminio LME vs. DXY}: Relación inversa entre el commodity y el índice global del dólar.}
    \label{fig:06}
\end{minipage}
\end{figure}

% ═════════════════════════════════════════════════════════════════════════════
% SECCIÓN 3: DINÁMICA DE LA INDUSTRIA GLOBAL DEL ALUMINIO
% ═════════════════════════════════════════════
\section{Dinámica de la Industria Global del Aluminio y Cointegración con el LME}

\subsection{Comportamiento Cíclico: LME, DXY y Cobertura Operativa}
Existe una correlación inversa entre el precio internacional de los metales básicos y la fortaleza del dólar global (DXY, \textbf{Figura \ref{fig:06}}). Aluar cuenta con una cobertura operativa natural: en fases contractivas del commodity acompañadas por apreciación del dólar, las correcciones del tipo de cambio real doméstico reducen los costos fijos en pesos (salarios y servicios), amortiguando la presión sobre el margen operativo.

\subsection{Análisis Econométrico: Test de Cointegración de Engle-Granger}
Se sometió la serie diaria de ALUA.BA (expresada en USD vía CCL, $N=2.355$ ruedas, período 2016--2026) y el futuro del LME a la prueba de cointegración de Engle-Granger. El estadístico $t = -2,29$ ($p = 0,38$, frente a un valor crítico del $-3,34$ al 5\%) no permite afirmar cointegración estricta en niveles. Sin embargo, la elasticidad contemporánea se sitúa en 0,66 ($R^2=0,12$), confirmando que Aluar responde de forma combinada tanto al ciclo internacional del metal como a las condiciones macroeconómicas locales.

\begin{figure}[H]
\centering
\begin{minipage}[t]{0.485\textwidth}
    \centering
    \includegraphics[width=\linewidth,height=4.6cm,keepaspectratio]{../03_Modelo_y_Codigo/figuras/figura_07.png}
    \caption{\textbf{Producción Mundial de Aluminio}: Participación de los principales países productores.}
    \label{fig:07}
\end{minipage}\hfill
\begin{minipage}[t]{0.485\textwidth}
    \centering
    \includegraphics[width=\linewidth,height=4.6cm,keepaspectratio]{../03_Modelo_y_Codigo/figuras/figura_08.png}
    \caption{\textbf{Capacidad y Autogeneración}: Comparativa de escala y matriz energética propia frente a pares.}
    \label{fig:08}
\end{minipage}
\end{figure}

\begin{figure}[H]
\centering
\begin{minipage}[t]{0.485\textwidth}
    \centering
    \includegraphics[width=\linewidth,height=4.6cm,keepaspectratio]{../03_Modelo_y_Codigo/figuras/figura_09.png}
    \caption{\textbf{Participación de Mercado Local}: Cuota del 60\% de Aluar frente al 40\% importado.}
    \label{fig:09}
\end{minipage}\hfill
\begin{minipage}[t]{0.485\textwidth}
    \centering
    \includegraphics[width=\linewidth,height=4.6cm,keepaspectratio]{../03_Modelo_y_Codigo/figuras/figura_10.png}
    \caption{\textbf{Matriz Renovable y Curva C1}: Posicionamiento en el primer cuartil de costos de fundición.}
    \label{fig:10}
\end{minipage}
\end{figure}

\subsection{Estructura del Mercado Internacional y Posición Local}
China produce el 56\% del aluminio primario mundial (\textbf{Figura \ref{fig:07}}), con una alta dependencia de generación eléctrica a carbón. Aluar, con una producción de 0,44 MM Tn anuales, abastece el \textbf{60\% del consumo interno argentino} (\textbf{Figura \ref{fig:09}}), fijando precios locales sobre la base de paridad de importación. Su matriz energética 78\% renovable la ubica en el percentil 29\% de la curva global de costos C1 (\textbf{Figura \ref{fig:10}}).

% ═════════════════════════════════════════════════════════════════════════════
% SECCIÓN 4: GOBERNANZA CORPORATIVA Y ESG
% ═════════════════════════════════════════════
\section{Gobernanza Corporativa, Marco Regulatorio y Factores ESG}

\subsection{Gobierno Corporativo y Transparencia (Ley 26.831 / CNV)}
Aluar cotiza en el régimen de oferta pública de la CNV (Resolución General 797/2019). El Directorio está encabezado por Javier Madanes Quintanilla. El Comité de Auditoría cuenta con mayoría de directores independientes según el Art. 109 de la Ley 26.831, dictaminando sobre el sistema de control interno y transacciones con sociedades vinculadas como Hidroeléctrica Futaleufú S.A. (60,2\%) y Transpa S.A. (20,4\%), celebradas a precios de mercado.

\begin{table}[H]
\centering
\caption{\textbf{Matriz de Evaluación de Factores Ambientales, Sociales y de Gobierno (ESG)}}
\scriptsize
\begin{tabularx}{\textwidth}{p{2.5cm}Xp{4.2cm}}
\toprule
\textbf{Dimensión} & \textbf{Indicadores Observables} & \textbf{Evaluación} \\
\midrule
\textbf{Ambiental (E)} & 78\% renovable propia. Huella $<$\,4 t CO\textsubscript{2}/t Al. Certificación ASI en proceso. & \textbf{Favorable.} Ventaja competitiva frente al mecanismo de ajuste en frontera por carbono (CBAM). \\
\textbf{Social (S)} & Principal empleador privado en Chubut ($>$\,6.200 puestos directos e indirectos). TRIR $<$\,1,0. & \textbf{Adecuada.} Estándares de seguridad industrial y baja tasa de siniestralidad. \\
\textbf{Gobernanza (G)} & Concentración del 45,98\% en el grupo de control. Auditoría externa por PwC y comités independientes. & \textbf{Estable.} Continuidad estratégica con flotante del 54,02\%. \\
\bottomrule
\end{tabularx}
\end{table}

\subsection{Impacto del Ajuste de Carbono en Frontera (CBAM de la Unión Europea)}
La reducida huella de carbono de Aluar ($<$\,4 t CO\textsubscript{2}/t Al) frente al promedio de la industria ($\sim$\,11 t CO\textsubscript{2}/t) le otorga una ventaja arancelaria neta estimada en \textbf{USD 80--120 por tonelada} bajo el esquema CBAM de la Unión Europea frente a productores alimentados con fuentes fósiles.

% ═════════════════════════════════════════════════════════════════════════
% SECCIÓN 5: VALUACIÓN RELATIVA Y DESAPALANCAMIENTO
% ═════════════════════════════════════════════
\section{Valuación Relativa y Dinámica de Desapalancamiento Post-CAPEX}

\subsection{Comparación de Múltiplos LTM y Forward}
Aluar cotiza a un múltiplo EV/EBITDA LTM de \textbf{13,4x} sobre el resultado deprimido de FY2025 (USD 163,3 MM). Al normalizar la proyección sobre el EBITDA estimado de FY2026E (\textbf{USD 391,2 MM}), el múltiplo implícito desciende a \textbf{5,2x EV/EBITDA Forward}, situando a la firma con un \textbf{descuento del 20\%--25\%} frente a comparables internacionales como Alcoa (8,1x) y Norsk Hydro (6,5x).

\begin{figure}[H]
\centering
\begin{minipage}[t]{0.485\textwidth}
    \centering
    \includegraphics[width=\linewidth,height=4.6cm,keepaspectratio]{../03_Modelo_y_Codigo/figuras/figura_11.png}
    \caption{\textbf{EV/EBITDA vs. Margen}: Ubicación de Aluar frente a pares internacionales (Alcoa, Hydro, Chalco).}
    \label{fig:11}
\end{minipage}\hfill
\begin{minipage}[t]{0.485\textwidth}
    \centering
    \includegraphics[width=\linewidth,height=4.6cm,keepaspectratio]{../03_Modelo_y_Codigo/figuras/figura_12.png}
    \caption{\textbf{EBITDA Histórico y Proyectado}: Recuperación operativa y desapalancamiento tras el pico de CAPEX.}
    \label{fig:12}
\end{minipage}
\end{figure}

\subsection{Trayectoria de Desapalancamiento Financiero}
El desembolso en bienes de uso durante FY2025 (USD 243,9 MM) elevó la Deuda Neta a USD 525,8 MM (3,22x EBITDA). Con la moderación de inversiones hacia niveles de mantenimiento ($\sim\text{USD } 90\text{--}103\text{ MM/año}$), el flujo libre de caja proyectado permite recomprimir la Deuda Neta / EBITDA a \textbf{$<1,0\text{x}$ hacia FY2028E} (\textbf{Figura \ref{fig:12}}).

\begin{table}[H]
\centering
\caption{\textbf{Múltiplos Financieros de Empresas Comparables del Sector (julio-2026)}}
\small
\begin{tabular}{llccc}
\toprule
\textbf{Compañía} & \textbf{Ticker} & \textbf{EV/EBITDA (x)} & \textbf{Margen EBITDA} & \textbf{Deuda Neta/EBITDA} \\
\midrule
Aluar S.A.I.C. (LTM / Fwd) & ALUA.BA & 13,4x / 5,2x & 14,9\% & 3,22x \\
Alcoa Corp. & AA & 8,1x & 9,2\% & 2,10x \\
Norsk Hydro & NHY.OL & 6,5x & 11,4\% & 1,40x \\
Chalco & 2600.HK & 5,8x & 8,8\% & 3,80x \\
\midrule
\textbf{Mediana de Pares} & & \textbf{6,5x} & \textbf{9,2\%} & \textbf{2,10x} \\
\bottomrule
\end{tabular}
\end{table}

% ═════════════════════════════════════════════════════════════════════════
% SECCIÓN 6: MODELO WACC Y FACTOR LAMBDA
% ═════════════════════════════════════════════
\section{Costo Promedio Ponderado del Capital (WACC) y Factor \texorpdfstring{$\lambda$}{lambda}}

\subsection{Formulación y Estimación del WACC}
El Costo Promedio Ponderado del Capital se calcula conforme a la formulación clásica:
\begin{equation}
    WACC = K_e \left(\frac{E}{V}\right) + K_d (1 - t) \left(\frac{D}{V}\right)
\end{equation}
Reemplazando los parámetros calibrados ($E/V = 67,29\%$, $D/V = 32,71\%$, $K_d = 3,80\%$, $t = 35\%$ y $K_e = 9,30\%$):
\begin{equation}
    WACC = 9,30\% \times 0,6729 + 3,80\% \times (1 - 0,35) \times 0,3271 = 6,26\% + 0,81\% = \mathbf{7,06\%}
\end{equation}

\subsection{Deducción del Beta en Cuatro Pasos (OLS \texorpdfstring{$\to$}{->} Blume \texorpdfstring{$\to$}{->} Hamada)}
\begin{enumerate}
    \item \textbf{Beta OLS Histórico:} Regresión lineal diaria de 10 años ($N=2.376$ ruedas) de ALUA (en USD CCL) contra el índice S\&P 500:
    \begin{equation}
        r_{ALUA, t} = \alpha + \beta_{OLS} r_{SP500, t} + \epsilon_t \implies \beta_{OLS} = 0,8420 \quad (SE = 0,0563, \; R^2 = 8,60\%)
    \end{equation}
    \item \textbf{Ajuste de Blume (1975):} Corrección por reversión a la media del riesgo sistemático:
    \begin{equation}
        \beta_{Blume} = \frac{2}{3} \beta_{OLS} + \frac{1}{3} (1,0) = \frac{2}{3} (0,8420) + \frac{1}{3} = \mathbf{0,8947}
    \end{equation}
    \item \textbf{Desapalancamiento por Hamada (1972):} Con el ratio $D/E$ histórico mediano ($D/E_{hist} = 0,5036$):
    \begin{equation}
        \beta_U = \frac{\beta_{Blume}}{1 + (1 - t) \left(\frac{D}{E}\right)_{hist}} = \frac{0,8947}{1 + (1 - 0,35) \times 0,5036} = \mathbf{0,6745}
    \end{equation}
    \item \textbf{Reapalancamiento al $D/E$ Objetivo ($D/E_{obj} = 0,4860$):}
    \begin{equation}
        \beta_L = \beta_U \left[ 1 + (1 - t) \left(\frac{D}{E}\right)_{obj} \right] = 0,6745 \times [1 + (1 - 0,35) \times 0,4860] = \mathbf{0,8876}
    \end{equation}
\end{enumerate}

\begin{figure}[H]
\centering
\begin{minipage}[t]{0.485\textwidth}
    \centering
    \includegraphics[width=\linewidth,height=4.6cm,keepaspectratio]{../03_Modelo_y_Codigo/figuras/figura_13.png}
    \caption{\textbf{Secuencia del Beta}: OLS (0,842) $\rightarrow$ Blume (0,895) $\rightarrow$ Desapalancado (0,675) $\rightarrow$ Hamada ($\beta_L = 0,8876$).}
    \label{fig:13}
\end{minipage}\hfill
\begin{minipage}[t]{0.485\textwidth}
    \centering
    \includegraphics[width=\linewidth,height=4.6cm,keepaspectratio]{../03_Modelo_y_Codigo/figuras/figura_14.png}
    \caption{\textbf{Composición del WACC}: Costo del equity $K_e = 9,30\%$, deuda post-tax $K_d = 2,47\%$, resultando en \textbf{7,06\%}.}
    \label{fig:14}
\end{minipage}
\end{figure}

\subsection{Fundamentación del Factor de Exposición Soberana \texorpdfstring{$\lambda = 0,20$}{lambda = 0,20}}
El costo del capital propio incorpora el ajuste por riesgo soberano (Damodaran; Dumrauf):
\begin{equation}
    K_e = R_f + \beta_L \times ERP + \lambda \times EMBI+ = 4,70\% + 0,8876 \times 4,18\% + 0,20 \times 4,41\% = \mathbf{9,30\%}
\end{equation}
El factor $\lambda = 0,20$ equivale a la porción de ventas colocadas en el mercado doméstico (20\%). Puesto que el 80\% restante se exporta a cotizaciones LME y los costos operativos fijos están denominados en pesos (salarios e insumos locales), los ajustes cambiarios licúan costos en dólares, actuando como cobertura natural frente al riesgo país.

\begin{table}[H]
\centering
\caption{\textbf{Sensibilidad del Costo de Capital y Precio Objetivo ante Variaciones en $\lambda$}}
\small
\begin{tabularx}{\textwidth}{Xcccc}
\toprule
\textbf{Parámetro de Riesgo País} & $\mathbf{\lambda = 0,20\text{ (Caso Base)}}$ & $\mathbf{\lambda = 0,50}$ & $\mathbf{\lambda = 0,75}$ & $\mathbf{\lambda = 1,00\text{ (Sin Cobertura)}}$ \\
\midrule
Prima por Riesgo País ($\lambda \times EMBI+$) & 0,88\% & 2,21\% & 3,31\% & 4,41\% \\
Costo del Equity ($K_e$) & 9,30\% & 10,63\% & 11,73\% & 12,83\% \\
\textbf{WACC Resultante (USD)} & \textbf{7,06\%} & \textbf{7,96\%} & \textbf{8,70\%} & \textbf{9,45\%} \\
\textbf{Precio Objetivo Base (ARS)} & \textbf{1.236,00} & \textbf{1.042,50} & \textbf{915,00} & \textbf{812,00} \\
Retorno Esperado vs. Spot (ARS 982,50) & +25,8\% & +6,1\% & -6,9\% & -17,4\% \\
\bottomrule
\end{tabularx}
\end{table}

% ═════════════════════════════════════════════════════════════════════════
% SECCIÓN 7: ANÁLISIS DE CREACIÓN DE VALOR (EVA)
% ═════════════════════════════════════════════
\section{Creación de Valor Económico (EVA) y Disciplina Terminal}

\subsection{Evolución del EVA y Efecto del Capital Invertido}
El Valor Económico Agregado ($EVA_t = (ROIC_t - WACC) \times NOA_{t-1}$) mide la generación de rentabilidad económica sobre los activos operativos netos. En FY2025 el EVA se situó en $-$USD 60,3 MM debido a la incorporación de USD 243,9 MM en activos eólicos al capital invertido (NOA $= \text{USD } 1.704,3\text{ MM}$) en forma previa a su plena entrada en régimen operativo.

\begin{table}[H]
\centering
\caption{\textbf{Trayectoria Histórica y Proyectada de Creación de Valor (EVA, USD MM)}}
\scriptsize
\begin{tabularx}{\textwidth}{Xcccccc|c}
\toprule
\textbf{Métrica de Valor} & \textbf{FY20} & \textbf{FY21} & \textbf{FY22} & \textbf{FY23} & \textbf{FY24} & \textbf{FY25} & \textbf{Terminal (2030E)} \\
\midrule
NOPAT & 31,3 & 39,9 & 105,4 & 40,3 & 99,2 & 60,0 & 135,5 \\
Capital Invertido (NOA) & 838,4 & 531,3 & 590,5 & 771,6 & 1.248,9 & 1.704,3 & 1.919,3 \\
ROIC (NOPAT / NOA) & 3,73\% & 7,50\% & 17,85\% & 5,22\% & 7,94\% & 3,52\% & \textbf{7,06\%} \\
WACC & 7,06\% & 7,06\% & 7,06\% & 7,06\% & 7,06\% & 7,06\% & \textbf{7,06\%} \\
Spread (ROIC - WACC) & -3,33\% & +0,44\% & +10,79\% & -1,84\% & +0,88\% & -3,54\% & \textbf{0,00\%} \\
\textbf{EVA (USD MM)} & \textbf{-27,9} & \textbf{2,3} & \textbf{63,7} & \textbf{-14,2} & \textbf{11,0} & \textbf{-60,3} & \textbf{0,0} \\
\bottomrule
\end{tabularx}
\end{table}

\subsection{Criterios Metodológicos: Caso Base Dumrauf vs. Cota de Reinversión Damodaran}
\begin{itemize}
    \item \textbf{Caso Base Oficial (Dumrauf, Cap. 14):} En estado estacionario a plena capacidad (460 kt), $CAPEX = D\&A$ y $\Delta NWC = 0$, por lo que $FCFF_{\text{terminal}} = NOPAT_{2030} = \text{USD } 135,5\text{ MM}$. El valor terminal descontado resulta $TV = \frac{135,5 \times 1,02}{0,0706 - 0,02} = \mathbf{USD\text{ }2.731,4\text{ MM}}$, determinando el precio objetivo oficial de \textbf{ARS 1.236,00}.
    \item \textbf{Cota de Reinversión (Damodaran, $ROIC \to WACC$):} Asumiendo una tasa de reinversión implícita $RR = g / WACC = 2,0\% / 7,06\% = 28,3\%$, el flujo terminal se ajusta a $NOPAT \times (1 - RR) = \text{USD } 97,1\text{ MM}$, arrojando un precio objetivo de \textbf{ARS 993,00}, que se adopta como cota conservadora de control.
\end{itemize}

% ═════════════════════════════════════════════════════════════════════════
% SECCIÓN 8: PUENTE DE VALUACIÓN
% ═════════════════════════════════════════════
\section{Puente de Valuación: Desglose de Enterprise Value a Equity Value}

\subsection{Composición del Enterprise Value}
El Enterprise Value estimado asciende a \textbf{USD 3.951 MM} (\textbf{Figura \ref{fig:15}}). Los flujos de fondos explícitos (2026E--2030E) aportan USD 680 MM (17,2\% del EV), mientras que el Valor Terminal representa USD 2.902 MM (73,4\% del EV), reflejando la vida útil indefinida de los activos de fundición. Deduciendo la Deuda Neta auditada de USD 525,8 MM y sumando activos no operativos, se obtiene un \textbf{Equity Value de USD 2.186 MM}, equivalente a \textbf{ARS 1.236,00 por acción} al tipo de cambio CCL de 1.584,25.

\begin{figure}[H]
\centering
\begin{minipage}[t]{0.485\textwidth}
    \centering
    \includegraphics[width=\linewidth,height=4.6cm,keepaspectratio]{../03_Modelo_y_Codigo/figuras/figura_15.png}
    \caption{\textbf{Puente de Valuación (Waterfall)}: De Enterprise Value (USD 3.951 MM) a Equity Value (USD 2.186 MM) y Target.}
    \label{fig:15}
\end{minipage}\hfill
\begin{minipage}[t]{0.485\textwidth}
    \centering
    \includegraphics[width=\linewidth,height=4.6cm,keepaspectratio]{../03_Modelo_y_Codigo/figuras/figura_16.png}
    \caption{\textbf{Capital de Trabajo Operativo}: Trayectoria de DIO, DSO, DPO y Ciclo de Conversión de Efectivo (CCC).}
    \label{fig:16}
\end{minipage}
\end{figure}

\subsection{Evolución del Ciclo de Conversión de Efectivo (CCC)}
El Ciclo de Conversión de Efectivo (\textbf{Figura \ref{fig:16}}) se comprime desde 79 días en FY2023 a \textbf{64 días proyectados}, producto de la normalización en días de inventario (DIO) y cobranza (DSO), favoreciendo la liquidez para el servicio de pasivos.

% ═════════════════════════════════════════════════════════════════════════
% SECCIÓN 9: ANÁLISIS ESTRATÉGICO DUPONT Y FRAMEWORKS
% ═════════════════════════════════════════════
\section{Análisis Estratégico DuPont y Fuerzas Competitivas}

\subsection{Descomposición DuPont de Cinco Factores (FY2020--FY2025)}
La descomposición DuPont extendida permite explicar la variabilidad del ROE histórico:
\begin{equation}
    ROE = \underbrace{\left(\frac{\text{RN}}{\text{EBT}}\right)}_{\text{Carga Impositiva}} \times \underbrace{\left(\frac{\text{EBT}}{\text{EBIT}}\right)}_{\text{Carga Financiera}} \times \underbrace{\left(\frac{\text{EBIT}}{\text{Ventas}}\right)}_{\text{Margen EBIT}} \times \underbrace{\left(\frac{\text{Ventas}}{\text{Activos}}\right)}_{\text{Rotación Activos}} \times \underbrace{\left(\frac{\text{Activos}}{\text{PN}}\right)}_{\text{Apalancamiento}}
\end{equation}
En FY2025, el factor impositivo se redujo a 15,67\% (alícuota efectiva del 84,3\% por el impacto del ajuste por inflación impositivo NIC 29), situando el ROE en 0,83\% aun cuando el margen operativo EBIT se ubicó en 8,44\%.

\begin{table}[H]
\centering
\caption{\textbf{Descomposición DuPont de 3 Factores (FY2020--FY2025)}}
\small
\begin{tabularx}{\textwidth}{Xcccccc}
\toprule
\textbf{Componente DuPont} & \textbf{FY2020} & \textbf{FY2021} & \textbf{FY2022} & \textbf{FY2023} & \textbf{FY2024} & \textbf{FY2025} \\
\midrule
Margen Neto ($RN / Ventas$) & -5,34\% & 5,48\% & 15,79\% & 19,48\% & 9,80\% & 0,84\% \\
Rotación Activos ($V / Activos$) & 0,78x & 0,74x & 0,79x & 0,71x & 0,63x & 0,55x \\
Apalancamiento ($Activos / PN$) & 2,27x & 2,00x & 1,82x & 1,68x & 1,74x & 1,78x \\
\midrule
\textbf{ROE (Producto)} & \textbf{-9,49\%} & \textbf{8,04\%} & \textbf{22,64\%} & \textbf{23,25\%} & \textbf{10,68\%} & \textbf{0,83\%} \\
\bottomrule
\end{tabularx}
\end{table}

\subsection{Las Cinco Fuerzas de Porter}
\begin{itemize}
    \item \textbf{Amenaza de Nuevos Entrantes (Muy Baja):} Elevada inversión de capital ($>$\,USD 2.500 MM para un complejo nuevo) y acceso restringido a fuentes hidroeléctricas dedicadas.
    \item \textbf{Poder de Negociación de Proveedores (Bajo a Medio):} Concentración en alúmina importada mitigada por contratos de suministro diversificados; insumo energético 96\% autogenerado.
    \item \textbf{Poder de Negociación de Clientes (Bajo):} Exportaciones colocadas a precios internacionales tomadores LME; cuota del 60\% en el mercado doméstico.
    \item \textbf{Amenaza de Sustitutos (Baja):} Propiedades mecánicas y reciclabilidad del aluminio difíciles de reemplazar en extrusión estructural y transporte.
    \item \textbf{Rivalidad Competitiva (Media):} Competencia internacional por costos; protección operativa en el primer cuartil C1.
\end{itemize}

\subsection{Marco Regulatorio: Impacto del Régimen RIGI}
El Régimen de Incentivos para Grandes Inversiones (RIGI) reduce la tasa corporativa del impuesto a las ganancias al 25\% para nuevos desarrollos eólicos y elimina aranceles de importación en bienes de capital, optimizando el rendimiento de expansiones futuras.

% ═════════════════════════════════════════════════════════════════════════
% SECCIÓN 10: ANÁLISIS FINANCIERO HISTÓRICO Y RATIOS
% ═════════════════════════════════════════════
\section{Análisis Financiero Histórico, Ratios Operativos y Cobertura}

\subsection{Desempeño Operativo y Cobertura de Intereses}
La serie histórica auditada demuestra la solidez operativa de Aluar en distintos momentos del ciclo macroeconómico. El ratio de cobertura de intereses ($EBIT/Intereses$) cerró en \textbf{3,53x en FY2025} tras el pico de endeudamiento financiero, con ratios de liquidez corriente por encima de 1,40x:

\begin{table}[H]
\centering
\caption{\textbf{Estados Financieros y Ratios Operativos Históricos de Aluar (USD MM)}}
\scriptsize
\begin{tabularx}{\textwidth}{Xcccccc}
\toprule
\textbf{Métrica / Ejercicio} & \textbf{FY2020} & \textbf{FY2021} & \textbf{FY2022} & \textbf{FY2023} & \textbf{FY2024} & \textbf{FY2025} \\
\midrule
Ventas Netas & 823,50 & 513,50 & 654,72 & 647,75 & 915,85 & 1.092,60 \\
EBITDA & 116,50 & 107,39 & 209,67 & 105,70 & 204,70 & 163,25 \\
Margen EBITDA & 14,15\% & 20,91\% & 32,02\% & 16,32\% & 22,35\% & 14,94\% \\
EBIT & 48,10 & 61,33 & 162,14 & 61,97 & 152,57 & 92,27 \\
Margen EBIT & 5,84\% & 11,94\% & 24,76\% & 9,57\% & 16,66\% & 8,44\% \\
Resultado Neto & -43,94 & 28,13 & 103,35 & 126,19 & 89,74 & 9,17 \\
Margen Neto & -5,34\% & 5,48\% & 15,79\% & 19,48\% & 9,80\% & 0,84\% \\
CAPEX & -103,01 & -9,43 & -11,11 & -65,89 & -25,59 & -243,93 \\
Deuda Neta & 351,30 & 137,44 & 57,33 & 169,88 & 211,23 & 525,76 \\
\midrule
\textbf{ROIC} & \textbf{3,73\%} & \textbf{7,50\%} & \textbf{17,85\%} & \textbf{5,22\%} & \textbf{7,94\%} & \textbf{3,52\%} \\
ROE & -9,49\% & 8,04\% & 22,64\% & 23,25\% & 10,68\% & 0,83\% \\
Liquidez Corriente & 1,89x & 3,78x & 2,92x & 2,92x & 4,93x & 2,44x \\
Deuda Neta / EBITDA & 3,02x & 1,28x & 0,27x & 1,61x & 1,03x & 3,22x \\
Cobertura Intereses ($EBIT/Int.$) & 1,81x & 4,12x & 107,80x & 7,69x & 8,62x & 3,53x \\
\bottomrule
\end{tabularx}
\end{table}

% ═════════════════════════════════════════════════════════════════════════
% SECCIÓN 11: VALUACIÓN FUNDAMENTAL (DCF)
% ═════════════════════════════════════════════
\section{Valuación Fundamental por Flujo de Fondos Descontados (DCF)}

\subsection{Estructura del Modelo DCF y Valor de la Opción Real PEAL V}
El valor patrimonial surge del descuento de flujos de fondos libres para la firma (FCFF) bajo la formulación canónica de Dumrauf:
\begin{equation}
    EV = \sum_{t=0}^{4} \frac{FCFF_{2026+t}}{(1 + WACC)^t} + \frac{TV}{(1 + WACC)^4} = \mathbf{USD\text{ }3.951,0\text{ MM}}
\end{equation}
\begin{itemize}
    \item \textbf{Precio Objetivo DCF Caso Base:} \textbf{ARS 1.236,00} por acción (USD 0,78, +25,8\% de retorno esperado).
    \item \textbf{Opción Real PEAL V (Longstaff-Schwartz):} Valúa la flexibilidad gerencial de diferir o ejecutar 312 MW eólicos adicionales mediante simulación por mínimos cuadrados (LSMC) sobre polinomios de Laguerre:
    \begin{equation}
        \hat{C}(S_t) = \mathbb{E}^{\mathbb{Q}}\left[ e^{-r \Delta t} V_{t+1} \mid S_t \right] = \sum_{j=0}^{2} a_j L_j(S_t)
    \end{equation}
    Aportando un valor aditivo de \textbf{+ARS 19,60 por acción} (+USD 0,01).
    \item \textbf{Precio Objetivo Integrado (Base + Opción Real):} \textbf{ARS 1.255,60} por acción (USD 0,79, \textbf{+27,8\% de retorno esperado}).
\end{itemize}

\subsection{Prueba de Sensibilidad ante Repunte del Riesgo Soberano a 2.400 pb}
En un escenario de tensión macroeconómica donde el EMBI+ repunta a 2.400 pb, el WACC asciende a 13,21\% (+615 pb) y el precio objetivo se retrae a ARS 237,00 (\textbf{Figura \ref{fig:17}}).

\begin{figure}[H]
\centering
\begin{minipage}[t]{0.485\textwidth}
    \centering
    \includegraphics[width=\linewidth,height=4.6cm,keepaspectratio]{../03_Modelo_y_Codigo/figuras/figura_17.png}
    \caption{\textbf{Sensibilidad al Riesgo Soberano}: Impacto sobre WACC y Target en escenario de EMBI+ a 2.400 pb.}
    \label{fig:17}
\end{minipage}\hfill
\begin{minipage}[t]{0.485\textwidth}
    \centering
    \includegraphics[width=\linewidth,height=4.6cm,keepaspectratio]{../03_Modelo_y_Codigo/figuras/figura_18.png}
    \caption{\textbf{Rango de Precios Intrínsecos}: Comparativa de estimaciones metodológicas vs. cotización Spot.}
    \label{fig:18}
\end{minipage}
\end{figure}

\subsection{Marco de Escenarios Ponderados (Pesimista / Base / Optimista)}
\begin{table}[H]
\centering
\caption{\textbf{Marco de Escenarios: Supuestos Operativos y Precios Objetivos Resultantes}}
\scriptsize
\begin{tabularx}{\textwidth}{p{2.6cm}XXX}
\toprule
\textbf{Variable} & \textbf{Pesimista (25\%)} & \textbf{Base (50\%)} & \textbf{Optimista (25\%)} \\
\midrule
Precio LME & Caída hacia costo marginal global ($\sim$USD 1.900/t). & Nivel medio de equilibrio OU (USD 2.450--2.814/t). & LME sostenido $>$\,USD 3.000/t por déficit de oferta. \\
Riesgo País (EMBI+) & Rebote soberano a $>800$ pb. & Estabilización a largo plazo ($\approx$441 pb). & Compresión continua por debajo de 350 pb. \\
WACC / g & 9,0\% / 1,5\% & 7,06\% / 2,0\% & 6,5\% / 2,2\% \\
\midrule
\textbf{Precio Objetivo (DCF)} & \textbf{ARS 781} (-20,5\%) & \textbf{ARS 1.236,00} (Caso Base) & \textbf{ARS 1.481} (+50,7\%) \\
\bottomrule
\end{tabularx}
\end{table}

% ═════════════════════════════════════════════════════════════════════════
% SECCIÓN 12: ANÁLISIS DE SENSIBILIDAD DE VALUACIÓN
% ═════════════════════════════════════════════
\section{Análisis de Sensibilidad Bidimensional y Matriz de Rangos}

\subsection{Sensibilidad Cruzada: WACC vs. Tasa de Crecimiento Terminal ($g$)}
La matriz bidimensional (\textbf{Figura \ref{fig:19}}) confirma que la recomendación COMPRAR se mantiene válida para cualquier tasa WACC inferior a 8,0\% con tasas de crecimiento terminal $g \ge 1,5\%$.

\begin{figure}[H]
\centering
\begin{minipage}[t]{0.485\textwidth}
    \centering
    \includegraphics[width=\linewidth,height=4.6cm,keepaspectratio]{../03_Modelo_y_Codigo/figuras/figura_19.png}
    \caption{\textbf{Matriz WACC vs. g}: Sensibilidad bidimensional del precio fundamental por acción en ARS.}
    \label{fig:19}
\end{minipage}\hfill
\begin{minipage}[t]{0.485\textwidth}
    \centering
    \includegraphics[width=\linewidth,height=4.6cm,keepaspectratio]{../03_Modelo_y_Codigo/figuras/figura_20.png}
    \caption{\textbf{Simulación Monte Carlo (20.000 iteraciones)}: Distribución empírica con innovaciones t-Student ($\nu=4,2$).}
    \label{fig:20}
\end{minipage}
\end{figure}

\begin{table}[H]
\centering
\caption{\textbf{Resumen de Precios Intrínsecos según Metodología (ARS/acción)}}
\small
\begin{tabularx}{\textwidth}{Xccc}
\toprule
\textbf{Metodología de Valuación} & \textbf{Mínimo} & \textbf{Central} & \textbf{Máximo} \\
\midrule
DCF $\lambda=1,0$ (Sin Cobertura de Riesgo País) & -- & 776 & -- \\
DCF con Tasa de Reinversión Damodaran & -- & 993 & -- \\
\textbf{DCF Caso Base Oficial (Dumrauf, $\lambda=0,20$)} & -- & \textbf{1.236} & -- \\
DCF Crecimiento Optimista ($g=2,5\%$) & -- & 1.371 & -- \\
Monte Carlo (Percentiles P5 / Mediana / P95) & 844 & 1.237 & 1.737 \\
\midrule
\textbf{Target Integrado (DCF + Opción Real PEAL V)} & -- & \textbf{1.255,60} & -- \\
Escenarios Ponderados (25/50/25) & -- & 1.184 & -- \\
Cotización Spot de Mercado & -- & 982,50 & -- \\
\bottomrule
\end{tabularx}
\end{table}

% ═════════════════════════════════════════════════════════════════════════
% SECCIÓN 13: SIMULACIÓN MONTE CARLO Y DCF INVERSO
% ═════════════════════════════════════════════
\section{Simulación Estocástica de Monte Carlo y DCF Inverso}

\subsection{Distribución Empírica con Innovaciones t-Student \texorpdfstring{($\nu=4,2$)}{(nu=4,2)}}
La simulación de Monte Carlo con 20.000 iteraciones e innovaciones con colas pesadas t-Student ($\nu = 4,2$, \textbf{Figura \ref{fig:20}}) aplica el factor de corrección $\sqrt{\frac{\nu-2}{\nu}}$:
\begin{equation}
    WACC_{sim} = WACC + \sigma_{WACC} \sqrt{\frac{\nu-2}{\nu}} \, T_\nu, \quad g_{sim} = g + \sigma_g \sqrt{\frac{\nu-2}{\nu}} \, T_\nu
\end{equation}
Arroja una mediana de \textbf{ARS 1.237} y un intervalo de confianza P5--P95 de \textbf{ARS 844 a ARS 1.737}, con un \textbf{85,2\% de probabilidad empírica de retorno positivo} frente a la cotización spot.

\subsection{Expectativas Implícitas: Tasa de Crecimiento del DCF Inverso \texorpdfstring{($g^* = 0,69\%$)}{(g* = 0,69\%)}}
Despejando la tasa de crecimiento terminal implícita en la cotización de mercado (\textbf{Figura \ref{fig:21}}), se comprueba que el precio actual descuenta un crecimiento a perpetuidad de apenas \textbf{$g^* = 0,69\%$}, muy por debajo del 2,0\% fundamental de largo plazo.

\begin{figure}[H]
\centering
\begin{minipage}[t]{0.485\textwidth}
    \centering
    \includegraphics[width=\linewidth,height=4.6cm,keepaspectratio]{../03_Modelo_y_Codigo/figuras/figura_21.png}
    \caption{\textbf{DCF Inverso}: Tasa g implícita descontada en el precio de mercado (0,69\%).}
    \label{fig:21}
\end{minipage}\hfill
\begin{minipage}[t]{0.485\textwidth}
    \centering
    \includegraphics[width=\linewidth,height=4.6cm,keepaspectratio]{../03_Modelo_y_Codigo/figuras/figura_22.png}
    \caption{\textbf{Convergencia de Múltiplos}: Normalización del múltiplo EV/EBITDA hacia 6,75x.}
    \label{fig:22}
\end{minipage}
\end{figure}

% ═════════════════════════════════════════════════════════════════════════
% SECCIÓN 14: VALUACIÓN RELATIVA SECTORIAL
% ═════════════════════════════════════════════
\section{Valuación Relativa Sectorial y Múltiplos Comparables LTM y Forward}

\subsection{Comparativa Internacional y Rendimiento del Flujo Libre de Caja}
El análisis comparativo sectorial confirma la divergencia entre la valuación basada en cifras LTM y la capacidad de generación normalizada. Mientras el múltiplo LTM de Aluar (13,3x) refleja la compresión de resultados de FY2025, el rendimiento proyectado del flujo libre de caja (\textit{FCF Yield 2027E} $>10\%$) supera ampliamente la mediana de la industria internacional (3,8\%):

\begin{table}[H]
\centering
\caption{\textbf{Múltiplos Financieros y Rendimientos del Sector Aluminio (julio-2026)}}
\scriptsize
\begin{tabularx}{\textwidth}{Xccccc}
\toprule
\textbf{Compañía / Ticker} & \textbf{EV/EBITDA LTM} & \textbf{Crec. Ventas LTM} & \textbf{Margen EBITDA LTM} & \textbf{ROIC LTM} & \textbf{FCF Yield 2026E/27E} \\
\midrule
ALUAR (ALUA.BA) & 13,3x & 8,0\% & 14,9\% & 4,3\% & -4,1\% / +12,8\% \\
Alcoa Corp. (AA) & 8,5x & 6,5\% & 12,0\% & 5,5\% & 4,8\% \\
Chalco (ACH) & 9,0x & 9,2\% & 15,1\% & 8,0\% & 5,2\% \\
Kaiser Aluminum (KALU) & 7,2x & 4,1\% & 9,8\% & 3,5\% & 3,0\% \\
Norsk Hydro (NHYDY) & 6,4x & 5,8\% & 13,5\% & 6,2\% & 4,1\% \\
Constellium (CSTM) & 5,4x & 3,2\% & 10,1\% & 5,0\% & 3,5\% \\
Rusal & 5,1x & -2,0\% & 8,4\% & 2,1\% & 1,5\% \\
\midrule
\textbf{Mediana de la Industria} & \textbf{7,2x} & \textbf{5,0\%} & \textbf{11,1\%} & \textbf{5,3\%} & \textbf{3,8\%} \\
\bottomrule
\end{tabularx}
\end{table}

% ═════════════════════════════════════════════════════════════════════════
% SECCIÓN 15: GESTIÓN CUANTITATIVA DE RIESGO
% ═════════════════════════════════════════════
\section{Gestión de Riesgos de Mercado y Asignación Óptima de Cartera}

\subsection{Criterio de Asignación de Kelly y Límite de Tolerancia por CVaR}
El criterio de crecimiento de Kelly arroja una fracción teórica óptima de:
\begin{equation}
    f^* = \frac{\mu - R_f}{\sigma^2} = \frac{24,63\% - 4,70\%}{(52,06\%)^2} = \mathbf{73,5\%} \text{ (Full-Kelly)} \implies Half-Kelly = \mathbf{36,8\%}
\end{equation}
Para la gestión prudencial de carteras se adopta un \textbf{límite de exposición del 20,0\%} (\textbf{Figura \ref{fig:23}}), acotado por el Conditional Value-at-Risk mensual al 95\% ($-32,80\%$). El Filtro de Kalman dinámico (\textbf{Figura \ref{fig:24}}) indica que el Beta de corto plazo converge a 0,764:
\begin{equation}
    y_t = \beta_t x_t + v_t, \quad \beta_t = \beta_{t-1} + w_t, \quad K_t = P_{t|t-1} x_t (x_t P_{t|t-1} x_t + R)^{-1}
\end{equation}
Esto demuestra que el beta Hamada de 0,888 utilizado en el WACC introduce un margen de seguridad conservador.

\begin{figure}[H]
\centering
\begin{minipage}[t]{0.485\textwidth}
    \centering
    \includegraphics[width=\linewidth,height=4.6cm,keepaspectratio]{../03_Modelo_y_Codigo/figuras/figura_23.png}
    \caption{\textbf{Criterio de Asignación de Kelly}: Full-Kelly (73,5\%), Half-Kelly (36,8\%) y Límite de Cartera (20,0\%).}
    \label{fig:23}
\end{minipage}\hfill
\begin{minipage}[t]{0.485\textwidth}
    \centering
    \includegraphics[width=\linewidth,height=4.6cm,keepaspectratio]{../03_Modelo_y_Codigo/figuras/figura_24.png}
    \caption{\textbf{Beta Dinámico por Filtro de Kalman}: Convergencia a 0,764 en el régimen macroeconómico reciente.}
    \label{fig:24}
\end{minipage}
\end{figure}

\subsection{Comportamiento Histórico de ALUA.BA en Períodos de Estrés}
\begin{table}[H]
\centering
\caption{\textbf{Desempeño Histórico de ALUA.BA (en USD) durante Episodios de Crisis}}
\small
\begin{tabularx}{\textwidth}{Xcccc}
\toprule
\textbf{Episodio de Estrés} & \textbf{Retorno Acumulado} & \textbf{Volatilidad Anualizada} & \textbf{Duración} & \textbf{Comportamiento} \\
\midrule
Crisis Covid-19 (Marzo 2020) & -44,73\% & 112,54\% & 50 ruedas & Shock de liquidez global \\
Tensión Cambiaria (Julio 2022) & +12,44\% & 45,91\% & 103 ruedas & Cobertura en activos dolarizados \\
Incertidumbre Electoral 2023 & +18,28\% & 105,36\% & 59 ruedas & Repreciación de capacidad industrial \\
Corrección LME (2022--2023) & +13,61\% & 44,40\% & 200 ruedas & Resiliencia por bajos costos operativos \\
\bottomrule
\end{tabularx}
\end{table}

\subsection{Estimación de Pérdidas Extremas por Teoría de Valores Extremos (EVT-GPD)}
Dado que los retornos de Aluar presentan colas pesadas (curtosis $= 7,82$), la aproximación normal subestima las pérdidas extremas. Se calibra un modelo Pareto Generalizado (GPD, $\hat{\xi} = 0,214 > 0$), capturando pérdidas de cola de tipo Fréchet:
\begin{equation}
    G_{\xi, \beta}(y) = 1 - \left(1 + \frac{\xi y}{\beta}\right)^{-1/\xi}
\end{equation}

\begin{table}[H]
\centering
\caption{\textbf{VaR y CVaR Diario por Metodología (retornos en USD, horizonte 1 día)}}
\small
\begin{tabularx}{\textwidth}{Xccc}
\toprule
\textbf{Metodología de Riesgo} & \textbf{VaR 90\%} & \textbf{VaR 95\%} & \textbf{VaR 99\%} \\
\midrule
Histórico & -3,47\% & -4,73\% & -8,20\% \\
Normal (paramétrico) & -4,12\% & -5,31\% & -7,54\% \\
t-Student ($\nu=4,46$) & -3,61\% & -4,99\% & -8,57\% \\
Cornish-Fisher & -2,70\% & -5,22\% & -13,41\% \\
CVaR / Expected Shortfall (paramétrico) & -5,67\% & -6,68\% & -8,65\% \\
CVaR Cornish-Fisher (Simulado) & -7,12\% & -10,49\% & -20,35\% \\
\bottomrule
\end{tabularx}
\end{table}

\begin{figure}[H]
\centering
\begin{minipage}[t]{0.485\textwidth}
    \centering
    \includegraphics[width=\linewidth,height=4.6cm,keepaspectratio]{../03_Modelo_y_Codigo/figuras/figura_25.png}
    \caption{\textbf{Valores Extremos (EVT-GPD)}: Estimación de pérdidas de cola con VaR 99\% (8,20\%) y ES 99\% (10,35\%).}
    \label{fig:25}
\end{minipage}\hfill
\begin{minipage}[t]{0.485\textwidth}
    \centering
    \includegraphics[width=\linewidth,height=4.6cm,keepaspectratio]{../03_Modelo_y_Codigo/figuras/figura_26.png}
    \caption{\textbf{Simulación LME (OU vs. MBG)}: Reversión a la media Ornstein-Uhlenbeck ($\theta=$ USD 2.814/t).}
    \label{fig:26}
\end{minipage}
\end{figure}

\subsection{Dinámica del Commodity: Proceso Ornstein-Uhlenbeck}
El aluminio LME exhibe reversión a la media hacia el costo marginal de producción de la industria (\textbf{Figura \ref{fig:26}}), modelado mediante la ecuación estocástica de Ornstein-Uhlenbeck:
\begin{equation}
    dS_t = \kappa (\theta - S_t) dt + \sigma S_t dW_t, \quad \kappa = 0,384, \; \theta = \text{USD 2.814/t}
\end{equation}

\begin{figure}[H]
\centering
\begin{minipage}[t]{0.485\textwidth}
    \centering
    \includegraphics[width=\linewidth,height=4.6cm,keepaspectratio]{../03_Modelo_y_Codigo/figuras/figura_27.png}
    \caption{\textbf{DCF Estocástico Continuo}: Densidad de probabilidad con 10.000 trayectorias simuladas.}
    \label{fig:27}
\end{minipage}\hfill
\begin{minipage}[t]{0.485\textwidth}
    \centering
    \includegraphics[width=\linewidth,height=4.6cm,keepaspectratio]{../03_Modelo_y_Codigo/figuras/figura_28.png}
    \caption{\textbf{Matriz de Riesgo}: Mapeo de Probabilidad vs. Impacto según estándares de análisis financiero.}
    \label{fig:28}
\end{minipage}
\end{figure}

% ═════════════════════════════════════════════════════════════════════════
% SECCIÓN 16: MODELIZACIÓN ESTOCÁSTICA Y ECONOMETRÍA
% ═════════════════════════════════════════════
\section{Modelización Estocástica Avanzada y Validación Econométrica}

\subsection{Validación de la Tasa de Descuento: AR(1) vs. Cox-Ingersoll-Ross (CIR)}
Para asegurar la no-negatividad matemática en las tasas de descuento del EMBI+ en simulaciones de colas extremas, se valida el modelo estocástico de Cox-Ingersoll-Ross (CIR, \textbf{Figura \ref{fig:29}}):
\begin{equation}
    dr_t = \kappa (\theta - r_t) dt + \sigma \sqrt{r_t} dW_t, \quad \text{Condición de Feller: } 2\kappa\theta > \sigma^2
\end{equation}
Con $\kappa = 0,7582$, $\theta = 0,1614$ y $\sigma = 0,085$, se verifica $2(0,7582)(0,1614) = 0,2448 > 0,0072 = \sigma^2$, garantizando que la tasa simulada nunca caiga por debajo de la tasa libre de riesgo de EE.UU.

\begin{figure}[H]
\centering
\begin{minipage}[t]{0.485\textwidth}
    \centering
    \includegraphics[width=\linewidth,height=4.6cm,keepaspectratio]{../03_Modelo_y_Codigo/figuras/figura_29.png}
    \caption{\textbf{Proceso CIR Soberano}: Simulación del EMBI+ respetando la no-negatividad estricta.}
    \label{fig:29}
\end{minipage}\hfill
\begin{minipage}[t]{0.485\textwidth}
    \centering
    \includegraphics[width=\linewidth,height=4.6cm,keepaspectratio]{../03_Modelo_y_Codigo/figuras/figura_30.png}
    \caption{\textbf{Beta Condicional GARCH(1,1)}: Variación temporal de la volatilidad y riesgo sistemático.}
    \label{fig:30}
\end{minipage}
\end{figure}

\begin{table}[H]
\centering
\caption{\textbf{Bondad de Ajuste para el Modelado del Riesgo Soberano}}
\small
\begin{tabularx}{\textwidth}{Xccc}
\toprule
\textbf{Modelo Estocástico} & \textbf{Log-Likelihood} & \textbf{AIC} & \textbf{BIC} \\
\midrule
AR(1) Lineal Base & 1.450,3 & -2.894,6 & -2.882,1 \\
\textbf{CIR (Estimación MLE)} & \textbf{1.682,7} & \textbf{-3.359,4} & \textbf{-3.346,9} \\
\bottomrule
\end{tabularx}
\end{table}

\subsection{Volatilidad Condicional y Beta GARCH(1,1)}
El modelo GARCH(1,1) (\textbf{Figura \ref{fig:30}}) describe la persistencia de varianza condicional:
\begin{equation}
    \sigma_t^2 = \omega + \alpha \epsilon_{t-1}^2 + \beta \sigma_{t-1}^2, \quad \alpha + \beta < 1
\end{equation}

\subsection{Dependencia Asimétrica en Caídas: Cópula de Clayton y Sensibilidad de Sobol}
La simulación multivariada adopta la Cópula de Clayton $C_\theta(u, v) = (u^{-\theta} + v^{-\theta} - 1)^{-1/\theta}$ frente a la Gaussiana ($\Delta AIC = -88,9$) debido a su capacidad de capturar dependencia de cola inferior ($\lambda_L = 2^{-1/\theta} = 0,38$), modelando adecuadamente la co-caída de los activos locales en episodios de pánico financiero.

\begin{figure}[H]
\centering
\begin{minipage}[t]{0.485\textwidth}
    \centering
    \includegraphics[width=\linewidth,height=4.6cm,keepaspectratio]{../03_Modelo_y_Codigo/figuras/figura_31.png}
    \caption{\textbf{Cópula de Clayton Asimétrica}: Modelado de co-caídas entre el activo y el mercado ($\lambda_L = 0,38$).}
    \label{fig:31}
\end{minipage}\hfill
\begin{minipage}[t]{0.485\textwidth}
    \begin{tcolorbox}[colback=white,colframe=gsnavy,arc=2pt,boxrule=0.8pt,top=3pt,bottom=3pt,left=3pt,right=3pt]
        \centering
        {\small\bfseries\color{gsnavy} Descomposición de Varianza (Sobol)}
        \vspace{2pt}
        \scriptsize
        \begin{tabularx}{\linewidth}{Xrr}
            \toprule
            \textbf{Variable} & \textbf{Índice $S_i$} & \textbf{Total $S_{Ti}$} \\
            \midrule
            Precio Aluminio LME & 0,421 & 0,510 \\
            Spread Soberano (EMBI) & 0,285 & 0,392 \\
            Costo Energía / Alúmina & 0,112 & 0,165 \\
            Tipo de Cambio Real & 0,084 & 0,140 \\
            \bottomrule
        \end{tabularx}
    \end{tcolorbox}
    \vspace{2pt}
    {\scriptsize Los índices de Sobol confirman que el 70,6\% del Enterprise Value depende conjuntamente del precio del LME y del spread soberano.}
\end{minipage}
\end{figure}

\subsection{Diagnóstico Econométrico y Pruebas Estadísticas}
\begin{table}[H]
\centering
\caption{\textbf{Resumen de Pruebas Estadísticas y Diagnóstico Econométrico}}
\scriptsize
\begin{tabularx}{\textwidth}{llccl}
\toprule
\textbf{Prueba} & \textbf{Variable / Proceso} & \textbf{Estadístico} & \textbf{P-Valor} & \textbf{Conclusión} \\
\midrule
Dickey-Fuller Aumentado (ADF) & Retornos Diarios ALUA (USD) & $t = -3,84$ & $p < 0,01$ & Serie estacionaria $I(0)$ \\
Phillips-Perron (PP) & Retornos Diarios ALUA (USD) & $Z_t = -4,12$ & $p < 0,01$ & Estacionariedad robusta a heterocedasticidad \\
Jarque-Bera & Normalidad de Retornos & $JB = 4.391,7$ & $p < 0,0001$ & Rechazo de normalidad (Curtosis = 7,82) \\
Kolmogorov-Smirnov & t-Student ($\nu=4,46$) vs. Normal & $D = 0,018$ & $p = 0,42$ & Ajuste con colas pesadas t-Student \\
Cópula de Clayton & Dependencia de Cola Inferior & $\Delta AIC = -88,9$ & $\lambda_L = 0,38$ & Dependencia asimétrica superior a Gaussiana \\
ARCH-LM (Lag 5) & Heterocedasticidad Condicional & $LM = 48,2$ & $p < 0,001$ & Presencia de volatilidad por conglomerados \\
Ljung-Box $Q(10)$ & Autocorrelación de Residuos & $Q = 8,45$ & $p = 0,58$ & Ausencia de correlación serial residual \\
Quandt-Andrews & Estabilidad Paramétrica & $F = 1,14$ & $p = 0,31$ & Parámetros estables en la muestra analizada \\
\bottomrule
\end{tabularx}
\end{table}

% ═════════════════════════════════════════════════════════════════════════
% SECCIÓN 17: CONCLUSIÓN
% ═════════════════════════════════════════════
\section{Dictamen Final y Conclusión del Modelo}

El modelo de valuación concluye una recomendación de \textbf{COMPRAR}, fundamentada en:
\begin{itemize}
    \item \textbf{Precio Objetivo Base (DCF Dumrauf):} \textbf{ARS 1.236,00} por acción (USD 0,78, upside de \textbf{+25,8\%}).
    \item \textbf{Precio Objetivo Integrado (Base + Opción Real PEAL V):} \textbf{ARS 1.255,60} por acción (USD 0,79, upside de \textbf{+27,8\%}).
    \item \textbf{Rango Probabilístico Monte Carlo (P5--P95):} ARS 844 a ARS 1.737 por acción, con un \textbf{85,2\% de probabilidad empírica de retorno positivo}.
    \item \textbf{Margen de Seguridad Fundamental:} La cotización actual descuenta una tasa terminal implícita contractiva ($g^* = 0,69\%$), ofreciendo un margen de seguridad de ARS 253,50 por acción frente al valor fundamental de régimen.
\end{itemize}

\clearpage

% ═════════════════════════════════════════════════════════════════════════
% APÉNDICE: ESTADOS FINANCIEROS AUDITADOS Y PROYECCIONES
% ═════════════════════════════════════════════
\section*{Apéndice: Estados Financieros Auditados y Proyecciones de Flujo de Fondos (USD MM)}
\addcontentsline{toc}{section}{Apéndice: Estados Financieros Auditados y Proyecciones de Flujo de Fondos}

{\small Cifras históricas anuales auditadas por Price Waterhouse \& Co. S.R.L. convertidas al CCL de cierre de cada ejercicio. Las proyecciones aplican la alícuota estatutaria del 35\% sobre el resultado operativo (Ley 20.628).}

\begin{table}[H]
\centering
\caption{\textbf{Conciliación de la Tasa Efectiva del Impuesto a las Ganancias FY2025}}
\small
\begin{tabularx}{\textwidth}{Xcc}
\toprule
\textbf{Componente Tributario} & \textbf{Monto (USD MM)} & \textbf{\% sobre EBT} \\
\midrule
Resultado Antes de Impuestos (EBT) & 58,5 & 100,0\% \\
Impuesto Teórico a Tasa Estatutaria (35\%) & (20,5) & 35,0\% \\
Efecto del Ajuste por Inflación Impositivo (NIC 29) & (23,1) & 39,5\% \\
Diferencias de Cambio no Deducibles y Otros & (5,7) & 9,8\% \\
\midrule
\textbf{Impuesto a las Ganancias Devengado en Estados Contables} & \textbf{(49,3)} & \textbf{84,3\%} \\
\bottomrule
\end{tabularx}
\end{table}

\begin{table}[H]
\centering
\caption{\textbf{Estado de Resultados Consolidado Auditado (FY2020--FY2025, USD MM)}}
\scriptsize
\begin{tabularx}{\textwidth}{Xrrrrrr}
\toprule
\textbf{Línea Contable} & \textbf{FY2020} & \textbf{FY2021} & \textbf{FY2022} & \textbf{FY2023} & \textbf{FY2024} & \textbf{FY2025} \\
\midrule
Ventas Netas (Revenue) & 823,5 & 513,5 & 654,7 & 647,8 & 915,9 & 1.092,6 \\
Costo de Ventas & (711,2) & (414,4) & (448,3) & (536,0) & (744,0) & (930,5) \\
\textbf{Ganancia Bruta} & \textbf{112,3} & \textbf{99,1} & \textbf{206,4} & \textbf{111,8} & \textbf{171,9} & \textbf{162,1} \\
Gastos de Administración y Comercialización & (63,8) & (37,8) & (44,1) & (50,1) & (70,1) & (100,8) \\
Otros Ingresos / Operativos Netos & (0,4) & 0,0 & (0,2) & 0,4 & 50,8 & 31,0 \\
\textbf{EBITDA} & \textbf{116,5} & \textbf{110,2} & \textbf{208,8} & \textbf{103,5} & \textbf{204,7} & \textbf{163,3} \\
Depreciación y Amortización (D\&A) & (68,4) & (48,9) & (46,7) & (41,5) & (52,1) & (71,0) \\
\textbf{EBIT (Resultado Operativo)} & \textbf{48,1} & \textbf{61,3} & \textbf{162,1} & \textbf{62,0} & \textbf{152,6} & \textbf{92,3} \\
Resultado Financiero y Asociadas & (84,5) & 14,7 & 36,2 & 88,4 & 3,9 & (33,8) \\
\textbf{Ganancia Antes de Impuestos (EBT)} & \textbf{(36,4)} & \textbf{76,0} & \textbf{198,3} & \textbf{150,3} & \textbf{156,5} & \textbf{58,5} \\
Impuesto a las Ganancias & (7,5) & (47,9) & (82,8) & (11,3) & (66,7) & (49,3) \\
\textbf{Resultado Neto del Ejercicio} & \textbf{(43,9)} & \textbf{28,1} & \textbf{115,5} & \textbf{139,0} & \textbf{89,7} & \textbf{9,2} \\
\textbf{NOPAT (Resultado Operativo Neto de Impuestos)} & \textbf{31,3} & \textbf{39,9} & \textbf{105,4} & \textbf{40,3} & \textbf{99,2} & \textbf{60,0} \\
\bottomrule
\end{tabularx}
\end{table}

\begin{table}[H]
\centering
\caption{\textbf{Estado de Situación Patrimonial / Balance Consolidado Auditado (FY2020--FY2025, USD MM)}}
\scriptsize
\begin{tabularx}{\textwidth}{Xrrrrrr}
\toprule
\textbf{Rubro Patrimonial} & \textbf{FY2020} & \textbf{FY2021} & \textbf{FY2022} & \textbf{FY2023} & \textbf{FY2024} & \textbf{FY2025} \\
\midrule
Activo Corriente & 454,6 & 318,1 & 478,9 & 490,1 & 902,9 & 1.028,1 \\
Activo No Corriente (inc. Bienes de Uso) & 597,4 & 379,9 & 353,6 & 422,3 & 556,6 & 940,9 \\
\textbf{Total Activo} & \textbf{1.052,0} & \textbf{698,0} & \textbf{832,6} & \textbf{912,4} & \textbf{1.459,5} & \textbf{1.969,0} \\
Pasivo Corriente & 240,1 & 84,2 & 163,8 & 167,8 & 183,2 & 420,5 \\
Pasivo No Corriente & 348,7 & 263,9 & 212,2 & 201,9 & 435,9 & 439,2 \\
\textbf{Total Pasivo} & \textbf{588,8} & \textbf{348,1} & \textbf{376,0} & \textbf{369,7} & \textbf{619,1} & \textbf{859,7} \\
Deuda Financiera Total & 375,2 & 181,4 & 134,0 & 228,9 & 408,5 & 594,9 \\
Caja, Bancos e Inversiones & 23,9 & 44,0 & 76,7 & 59,0 & 197,3 & 69,2 \\
\textbf{Deuda Financiera Neta} & \textbf{351,3} & \textbf{137,4} & \textbf{57,3} & \textbf{169,9} & \textbf{211,2} & \textbf{525,8} \\
\textbf{Patrimonio Neto} & \textbf{463,2} & \textbf{349,9} & \textbf{456,5} & \textbf{542,7} & \textbf{840,4} & \textbf{1.109,3} \\
\textbf{Capital Invertido (NOA)} & \textbf{838,4} & \textbf{531,3} & \textbf{590,5} & \textbf{771,6} & \textbf{1.248,9} & \textbf{1.704,3} \\
\bottomrule
\end{tabularx}
\end{table}

\begin{table}[H]
\centering
\caption{\textbf{Estado de Flujos de Efectivo Consolidado Auditado (FY2020--FY2025, USD MM)}}
\scriptsize
\begin{tabularx}{\textwidth}{Xrrrrrr}
\toprule
\textbf{Flujo de Caja} & \textbf{FY2020} & \textbf{FY2021} & \textbf{FY2022} & \textbf{FY2023} & \textbf{FY2024} & \textbf{FY2025} \\
\midrule
Flujo de Caja Operativo (FCO) & 245,7 & 102,4 & 89,2 & 113,6 & 79,2 & 30,7 \\
Inversiones en Bienes de Uso (CAPEX) & (103,0) & (9,4) & (11,1) & (65,9) & (25,6) & (243,9) \\
Flujo de Caja de Inversión (FCI) & (102,1) & (9,4) & (11,1) & (65,9) & (25,6) & (244,0) \\
Flujo de Caja de Financiación (FCF) & (197,8) & (45,1) & (34,4) & (70,8) & 91,1 & 20,3 \\
\textbf{Variación Neta de Efectivo} & \textbf{(54,2)} & \textbf{47,8} & \textbf{43,8} & \textbf{(23,1)} & \textbf{144,7} & \textbf{(193,0)} \\
\bottomrule
\end{tabularx}
\end{table}

\begin{table}[H]
\centering
\caption{\textbf{Proyecciones del Free Cash Flow to the Firm (FCFF, FY2026E--FY2030E, USD MM)}}
\scriptsize
\begin{tabularx}{\textwidth}{Xrrrrr}
\toprule
\textbf{Cuenta / Ejercicio} & \textbf{2026E} & \textbf{2027E} & \textbf{2028E} & \textbf{2029E} & \textbf{2030E} \\
\midrule
Ventas Netas (Revenue) & 1.591,4 & 1.543,0 & 1.494,6 & 1.446,3 & 1.397,9 \\
EBITDA & 391,2 & 369,3 & 347,4 & 325,5 & 303,6 \\
(-) Depreciación y Amortización & (108,4) & (105,1) & (101,8) & (98,5) & (95,2) \\
\textbf{EBIT} & \textbf{282,8} & \textbf{264,2} & \textbf{245,6} & \textbf{227,0} & \textbf{208,4} \\
(-) Impuestos Operativos (35\% Estatutaria) & (99,0) & (92,5) & (86,0) & (79,4) & (72,9) \\
\textbf{NOPAT} & \textbf{183,8} & \textbf{171,7} & \textbf{159,6} & \textbf{147,6} & \textbf{135,5} \\
(+) D\&A (Reversión) & 108,4 & 105,1 & 101,8 & 98,5 & 95,2 \\
(-) CAPEX de Mantenimiento & (103,2) & (100,0) & (96,9) & (93,8) & (90,6) \\
(-) Variación Capital de Trabajo ($\Delta$NWC) & (246,5) & 23,9 & 23,9 & 23,9 & 23,9 \\
\midrule
\textbf{FCFF Explícito} & \textbf{-57,5} & \textbf{200,7} & \textbf{188,4} & \textbf{176,2} & \textbf{164,0} \\
\bottomrule
\end{tabularx}
\end{table}

\noindent
{\small \textbf{Reconciliación FCFF Explícito 2030E vs. Terminal Normalizado.} FCFF Explícito 2030E: USD 164,0 MM (con CAPEX de mantenimiento de USD 90,6 MM y liberación de NWC de USD 23,9 MM). FCFF Terminal Normalizado ($g = 2,0\%$): USD 135,5 MM (asume estado estacionario donde CAPEX = D\&A y $\Delta$NWC = \$0, coincidiendo con el NOPAT de 2030E según la formulación canónica de Dumrauf Cap. 14).\par}

\clearpage

% ═════════════════════════════════════════════════════════════════════════
% REFERENCIAS BIBLIOGRÁFICAS Y AVISO LEGAL
% ═════════════════════════════════════════════
\section*{Referencias Bibliográficas y Fuentes de Información}
\addcontentsline{toc}{section}{Referencias Bibliográficas y Fuentes de Información}
\begin{enumerate}
    \small
    \item Damodaran, Aswath (2012). \textit{Investment Valuation: Tools and Techniques for Determining the Value of Any Asset}, 3rd Ed., John Wiley \& Sons.
    \item Dumrauf, Guillermo L. (2013). \textit{Finanzas Corporativas: Un Enfoque Latinoamericano}, 3ª Ed., Alfaomega Grupo Editor.
    \item Ross, Stephen A., Westerfield, Randolph W., Jaffe, Jeffrey (2019). \textit{Corporate Finance}, 12th Ed., McGraw-Hill Education.
    \item Alexander, Gordon J., Sharpe, William F., Bailey, Jeffery V. (2001). \textit{Fundamentos de Inversiones}, Pearson Educación / Prentice Hall.
    \item Black, Fischer, Scholes, Myron (1973). ``The Pricing of Options and Corporate Liabilities'', \textit{Journal of Political Economy}, 81(3), pp. 637--654.
    \item Longstaff, Francis A., Schwartz, Eduardo S. (2001). ``Valuing American Options by Simulation: A Simple Least-Squares Approach'', \textit{Review of Financial Studies}, 14(1), pp. 113--147.
    \item Blume, Marshall E. (1975). ``Betas and Their Regression Tendencies'', \textit{Journal of Finance}, 30(3), pp. 785--795.
    \item Hamada, Robert S. (1972). ``The Effect of the Firm's Capital Structure on the Systematic Risk of Common Stocks'', \textit{Journal of Finance}, 27(2), pp. 435--452.
    \item Sklar, Abe (1959). ``Fonctions de Répartition à n Dimensions et Leurs Marges'', \textit{Publications de l'Institut de Statistique de l'Université de Paris}, 8, pp. 229--231.
    \item Rockafellar, R. Tyrrell, Uryasev, Stanislav (2000). ``Optimization of Conditional Value-at-Risk'', \textit{Journal of Risk}, 2(3), pp. 21--42.
    \item Jaschke, Stefan R. (2001). ``The Cornish-Fisher-Expansion in the Context of Delta-Gamma-Normal Approximations'', \textit{Journal of Risk}, 4(4), pp. 33--52.
    \item Engle, Robert F., Granger, Clive W.J. (1987). ``Co-integration and Error Correction: Representation, Estimation and Testing'', \textit{Econometrica}, 55(2), pp. 251--276.
    \item Aluar Aluminio Argentino S.A.I.C. (2020–2025). \textit{Estados Financieros Consolidados Auditados y Memorias Anuales}, Comisión Nacional de Valores (CNV) \& Bolsas y Mercados Argentinos (BYMA).
\end{enumerate}

\vspace{0.4cm}
\noindent
\begin{tcolorbox}[colback=lightgray!30, colframe=gsnavy, boxrule=0.8pt, arc=3pt, title={\small\sffamily \textbf{Aviso Legal y Certificación del Analista (Ley 26.831)}}, coltitle=white]
\scriptsize
Este documento fue elaborado con fines exclusivamente de investigación académica y análisis financiero por Federico Agustín Chillón (FCE, Universidad Nacional de Cuyo). El presente reporte \textbf{no constituye oferta pública, asesoramiento financiero personalizado, recomendación vinculante de compra o venta, ni intermediación bursátil} en los términos de la Ley de Mercado de Capitales N.\textsuperscript{o}~26.831 y normas concordantes de la Comisión Nacional de Valores (CNV). El autor certifica que las estimaciones y modelizaciones reflejan un análisis técnico independiente fundamentado en estados contables públicos auditados por Price Waterhouse \& Co. S.R.L., series de mercado de BYMA, CBOE, S\&P Dow Jones y el London Metal Exchange (LME). Cualquier decisión de inversión que un tercero adopte en base a este informe es de su exclusiva responsabilidad.
\end{tcolorbox}

\end{document}
